\chapter{Introduction \label{chap1}}

Scattering, absorption and emission of light are central problems in science and engineering. Examples of these phenomena can be found, for instance, in earth's atmosphere, such as the colour of the sky or the clouds, in colloidal suspensions such as milk, or in biological tissues such as eyes. Because light scattering is strongly dependent of the particles size, shape and refractive index, scattering measurements are a crucial tool for astrophysics, biology, and engineering.\\
In the last decades, with the advent of nanophotonics, this topic has gain popularity. Nanophotonics or nano-optics is the branch of optics devoted to the study of light interaction with matter at the nanometer scale. The typical wavelength of the electromagnetic radiation lies between the ultraviolet ($300nm$)  and the near-infrared ($1.4\mu m$). This range of the electromagnetic spectrum is of special interest because the photon's energy lies in the range of the electronic transitions of matter. The human vision is adapted to detect radiation at these wavelengths, thus resulting in the diversity of colours. At these dimensions the electromagnetic radiation is well described by the wave picture and the use of classical theory of fields based on Maxwell's equations is suitable. 
%However, the small dimensions of matter require a quantum description of its properties. So, in nanophonics, is usual the use of a semiclassical theory, where classical theory of fields is combined with the quantum picture of matter.\\

%The aim of this thesis is the study of light interaction with nonreciprocal media and disordered systems, both at nanometric scale. In the first part we study emission and transport phenomena in magnetoplasmonic nanostructures, where the influence of an external magnetic field in the global behaviour of the system is analysed. The second part is devoted to the study of emission phenomenon in random correlated systems. 

\section{Light scattering with matter}

The propagation of light in a homogeneous and linear medium is done in straight paths. When inhomogeneities are present in the system, like particles or interfaces, light can be deflected from its straight path and consequently the direction of propagation changes.\\ 
Analysing the scattering systems, two different regimes can be identified: single and multiple scattering. In the single scattering regime, the field resulting from the interaction is described by the sum of the scattered fields (directly scattered, only once) from each of the scatterers in the system. Multiple scattering enters when the previous picture breaks down. In this case, the contributions of "partial waves" that have interacted with several scatterers are relevant. In the multiple scattering there is interplay between density, scattering cross section and dimensionality of the problem.\\
%When scattering occurs due to the presence of only one scatterer, or if in the system are present many scatterers but the scattered light don't interfere again with other scatterers, the system is in the single scattering regime; when light interacts many times with scatterers is called multiple scattering.\\
The treatment of light interaction requires different descriptions, depending of object scale under study. For objects with dimensions much bigger than the wavelength, the interaction can be described by geometrical optics, which considers light propagation in terms of rays.
This simplification is an excellent approximation to address many problems at the macroscale, like reflection or refraction phenomena including optical systems. It is used to describe the geometrical aspects
of imaging, including optical aberrations. However, it fails when  diffraction or interference phenomena are relevant in the system.\\
For objects with dimensions similar to the wavelength, it is common to consider the formal solution of Maxwell's equations with appropriate boundary conditions \cite{hulst1957light, bohren2008absorption}.
This approach is very useful due to the existence of exact solutions, although can only be applied to a small number of simple shapes such as spheres (usually called Mie solution to Maxwell's equations \cite{mie1908beitrage}) cylinders, spheroids, etc. It is appropriate to describe droplets in emulsions like milk, water droplets in the atmosphere, biological cells or interstellar grains.\\
When light is scattered by objects much smaller than the wavelength, like individual atoms, molecules or small particles, it is called Rayleigh scattering \cite{strutt1871lviii, rayleigh1899xxxiv}. Lord Rayleigh, in the end of the XIX century, observed that the scattering cross section of a small particle was dependent of the wavelength, $\propto\lambda^{-4}$. This dependence means that for short wavelengths the scattering is stronger than for long wavelengths. The typical example is the scattering of sunlight in the atmosphere, where the blue radiation (shorter in wavelength) is strongly scattered than red radiation (longer in wavelength). This causes the blue sky during the day and the red colour at sunset. When light is incident over one of these small particles, the oscillating electric field of the wave excites the charges within the particle, causing them to move at the same frequency. A common description of these systems is to consider each particle as a simple radiating dipole, characterised by its polarisability.\\
In the multiple scattering regime, like said above, light interacts many times with scatterers. In order to characterise the scattering strength of a system, the scattering mean free path $l_{s}$ is frequently used. This fundamental quantity determines the average distance travelled by a wave between two scattering events. The mean free path is given by $l_{s} = \left( \rho \sigma_{s} \right)^{-1}$, where $\rho$ is the density of the scatterers and $\sigma_{s}$ is the average scattering cross section, {\it{i.e.}}, the amount of light removed from the incident beam by scattering. Phenomenologically, the multiple scattering is the responsible for the turbid or opaque look of some media, like in a vapour cloud.\\
% In this scattering regime, the scattering mean free path is much smaller than the system dimensions, $l_{s}\ll L$. The multiple scattering is the responsible for the turbid or opaque look of some media, like in a vapour cloud.\\
% where the propagation of light is constantly changing, making part of the radiation towards back to the source.\\
The scattering phenomenon can be classified as elastic or inelastic. The first type, elastic scattering, the interacting light preserves the frequency but alters the propagation direction. The second type, inelastic scattering, alters both the wave frequency and the propagation direction. In this thesis, the systems under study obey to the first type.\\
In complex systems, when the distance between scatterers is sufficiently large or the scattering cross section is too small, the single scattering approximation can be used. However, for reduced distances or big scattering cross sections, the system gets in the multiple scattering regime and the single scattering is no longer valid. Depending on the scatterers positioning, we can categorise the systems in two extreme situations: the crystalline structure on one side and the disordered media on the other.\\
The canonical example of ordered structures are photonic crystals \cite{soukoulis2001photonic, galisteo2011self}. Very similar to the ionic lattices in solids, photonic crystals are periodic nanostructures that affect the propagation of light due to the regular variation of the refraction index. Proposed by E. Yablonovitch \cite{yablonovitch1987inhibited} and S. John \cite{john1987strong}, the multiple scattering due to the periodic variation of the refraction index causes a splitting of the bands at the edges of the Brillouin zone, called stop gap. When a stop gap is independent of the direction is refereed to as full band gap. The propagation of light through a photonic crystal is forbidden for light at the same energy as the band gap. At optical wavelengths, the existence of a photonic crystal full band gap was experimentally demonstrated by Blanco \textit{et al.} \cite{blanco2000large}.\\
Disordered media, in this context, implies a random distribution of the scatterers. The random distribution of the system can result in very different effects when light propagates through them. In next section the subject is discussed in detail.

%To evaluate the scattering strength of a system is common the use of the Ioffe-Regel parameter \cite{ioffe1960non} $\xi_{\textnormal{IR}}=k l_{s}$, where $k$ is the wavenumber in the medium and $l_{s}$ the scattering mean free path that determine the average distance travelled by the wave between two collisions or scattering events. The mean free path is defined by $l_{s} = \left( \rho \sigma_{s} \right)^{-1}$, where $\rho$ is the density of the scatterers and $\sigma_{s}$ is the average scattering cross section, {\it{i.e.}}, the amount of light removed from the incident beam by scattering.
%In a very diluted medium, where the scatterers density is low, and/or the scattering cross section is small, the medium is a weakly-scattering medium and the $\xi_{\textnormal{IR}}\llap 1$. In a compact media, and/or when the scattering cross section is big, the scattering can be strong enough to localise light. In this conditions $\xi_{\textnormal{IR}}\simeq 1$, this being the criterion of localisation \cite{ioffe1960non}.\\


\section{Light scattering in complex media \label{sec:Light_comp_media}}

In nature, the vast majority of interactions between light and matter occur with complex media, {\it{i.e.}}, media that is constituted by many parts. We can find examples almost everywhere, from fruits to vervet monkey skin \cite{vignolini2012pointillist,vukusic2007brilliant,Burresi:cq,price1975control}.\\
In the past, periodicity and spatial order was considered crucial to photonics, where many advances have been made with the goal of minimise and eliminate disorder. The control and reproducibility of ordered structures make them attractive to applications. A clear advantage was the possibility of reproduce the same behaviour when light interacts with them. In particular, photonic crystals have experienced a remarkable growth \cite{galisteo2011self}. However, in recent years researchers understood that disordered structures may be used as a powerful ally to control and manipulate light. \\
When light propagates through a disordered medium, like a white marble wall, it is multiple scattered by the inhomogeneities, resulting in a scrambling of the light. Although the incident light is very different from the transmitted through the medium, if the involved scattering process is elastic, the information is preserved \cite{wiersma1997localization}. The complex output or the transmitted intensity pattern is called speckle pattern and is the result of the interference between randomly scattered waves. An important research line is the identification of the descrambling key to reconstruct the input wave \cite{mosk2012controlling, bertolotti2012non,popoff2010measuring}, in particular through the so-called wavefront-shaping technique. Unfortunately, when the light source is embedded in a random structure this reconstruction seems more difficult. The solution of this problem is a topic of maximum importance due to the possible applications in biological imaging \cite{gibson2005recent}. \\
%Here, the constructive interference between waves coun- terpropagating along a given multiple scattering path enhances the average diffuse intensity reflected from an ensemble of random samples in a narrow angular range around the backscattering direction [6]. Under optimal experimental conditions allowing to apply the reciprocity theorem [7], this two-wave interference enhances the in- tensity exactly by a factor of 2 [8]
The interference and propagation of light in random materials can originate effects like weak localisation and coherent backscattering. In these phenomena waves follow the same path but in opposite directions, where always interference constructively and give rise to a robust wave phenomenon. Let us suppose that a light wave is incident over a finite random structure, and we collect the light in the backscattering direction. In this system, a large number of light paths may exist from the source to the detector, each one following a random trajectory through the media and then being backscattered out. In this case, due to the reciprocity, each trajectory can be followed in two directions, resulting in an enhancement in the backscattering direction. Surrounding this direction the enhancement gradually disappear, generating a cone. The phenomenon was experimentally measured by Albada \textit{et al.} \cite{VanAlbada:1985bx} and observed in biological tissues, in particular breast and lung tissues \cite{Yoo:1990vb}.\\
Another phenomenon with origin in the interference and propagation of light in random materials is the Anderson localisation \cite{Anderson:1958fz}. Proposed by Anderson to explain the metal-insulator transitions in electron transport, it is a wave phenomenon where disorder causes the interruption of transport. Equivalent phenomenon is present in optics, where modes with a high level of spatial confinement are created by multiple scattering, in which light remains in the system for a long time \cite{john1984electromagnetic, anderson1985question}. 
%Similarly to the weak localisation, Anderson localisation can be seen as wave paths that form closed loops, confining light in space.
In optics, efforts have been made in order to experimentally demonstrate this phenomenon \cite{wiersma1997localization,storzer2006observation}, but in three dimensional systems it is still an open problem \cite{wiersma2013disordered}.

Although most of the studies focus on perfectly ordered or disordered systems, between these two limits there is a region with great potential to explore. When the positions of the scatterers are correlated, the optical properties of the system can dramatically change. 
Random correlated structures, can display wide-angle reflection and broad spectral features. Some experiments show that reflection and transmission properties are governed by an interplay between Mie scattering and local order, by modification of the single scattering cross section \cite{reufer2007transport,Rojas_2004, garcia2007photonic}. In the past, relevant studies about conductivity in liquid metals \cite{Ashcroft1966} or light scattering in the cornea \cite{hart1969light} emphasised the importance of the correlated systems.\\
 In nanophotonics, these configurations can be achieved by introducing disorder in ordered systems or vice-versa. The standard approach considers a periodic system in which disorder is introduced by randomly removing elements. The consequent cavities, which are randomly distributed through the system, can couple. 
 %While the randomness of the system leads to a strong multiple scattering, the low density of optical modes within in a certain frequency window is assured by the long-range order. 
 S. John suggested that this combination may form an efficient trapping mechanism to localise light and obtain strong absorption \cite{john1983wave,john1984electromagnetic}.\\
 An alternative approach is based on the modification of a random structure by introducing correlations between positions of scattering elements. This can be done by considering charged particles, where the long-range electrostatic repulsion creates a short-range positional order \cite{john1983wave, Rojas_2004}. This type of manipulation allows the design of structures with a high control of the scattering strength.\\
The quest for the control of correlations was boosted in recent years by the emergence of the hyper-uniform structures \cite{florescu2009designer, wiersma2013disordered}. The design of these structures is based on a mathematical model \cite{torquato2003local} where the density fluctuations of hyper-uniform structures are reduced compared with those of a random system. Photonic crystals, quasicrystals and other disordered systems are included in the set of hyper-uniform structures. Recently Haberko \textit{et al.} \cite{haberko2013fabrication} fabricated the first three-dimensional photonic structures, using direct laser writing, based on a hyper-uniform structure. It was shown that these random structures exhibit a photonic bandgap \cite{florescu2009designer,muller2014silicon}.\\
The illumination of a disordered system with coherent light, as mentioned above, generates a speckle pattern. This pattern has a spatial structure that is usually characterised by the intensity spatial correlation function $ \left\langle I\left(\bold{r}\right)I \left(\bold{r}'\right) \right\rangle$. In transmission, short-range and long-range contributions can be identified in the intensity correlation function, that is function of the sum of three different terms denoted by $C_{1}$ (short-range), $C_{2}$ (intermediate-range) and $C_{3}$ (long-range) \cite{feng1988correlations}.\\
In systems where the source is located inside of a medium, the situation changes and a new type of contribution appears, the $C_{0}$ \cite{shapiro1999new}. This term  has attracted interest because it is very sensitive to the environment where the source is embedded \cite{skipetrov2000nonuniversal}. Also it has been shown that the $C_0$ correlations are equal to the fluctuations of the local density of states (LDOS) at the point source emitter \cite{skipetrov2006dynamics}. Recently was presented a non-diagrammatic approach of $C_{0}$, based on the fluctuations of LDOS and energy conservation \cite{caze2010near}.\\
In nanophotonics (although is a general result), the LDOS fluctuations can be obtained by measuring the decay rate fluctuations of a single emitter \cite{birowosuto2010observation,Sapienza_2011}. The modification of the behaviour of an emitter by the presence of the electromagnetic environment is known as Purcell \cite{Purcell_1946} effect. Experimentally, this phenomenon has been observed in emitters placed close to planar interfaces \cite{Drexhage_1966}, cavities \citep{Goy_1983}, two \cite{Kress_2005, Englund_2005} and three dimensional photonic crystals \cite{Lodahl_2004} and more recently in plasmonic \cite{Carminati_2006, Castanie_2012_distance} and magnetoplasmonic structures \cite{Nikolova_2013}.\\
The interest for the study of disordered systems with emission experiments dates back to the 90's, where Martorell \textit{et al.} \cite{Martorell_1991} studied the lifetime modification of a fluorescent molecule in colloidal suspensions. In this experiment, the concentration of particles in the colloid was varied to control the correlations between scatterers, and consequently a modification of the spontaneous emission was measured. This work directly relates the importance of correlations in the emission phenomenon. 
Recently, Krachmalnicoff \textit{et al.} \cite{krachmalnicoff2010fluctuations} have measured the statistical distribution of the LDOS on disordered semicontinuous gold films. The considered metallic films exhibited ''hot-spots'', \textit{i.e.}, high values of the electric field confined in deeply subwavelength areas. It was shown that LDOS fluctuations were maximum in the regime where the ''hot-spots'' were expected to dominate.
Birowosuto \textit{et al.} \cite{birowosuto2010observation} observed, in three dimensional random photonic media, fluctuations of the LDOS by measuring spontaneous emission of individual fluorescent nanospheres embedded in the structures. Sapienza \textit{et al.} \cite{Sapienza_2011} in a similar experiment, studied the Purcell effect for nano-emitters buried in a three dimensional medium. Making a massive number of measurements, they found that the statistics of Purcell factors presented non-Gaussian long-tailed distributions, that were explained by considering the full near-field multiple scattering.\\
These experiments highlight the importance of the LDOS distribution in applications in sensing or imaging of microscopic structures of complex media. However, it is not clear how sensitive the LDOS is to subtle modifications of the environment, like in phase transitions, or if can be used to monitor structural dynamic processes.

%A better understanding of these systems certainly implies a deeper understanding of the importance of disorder in photonics.

\section{Light scattering in nonreciprocal media}

Non-reciprocity refers to the impossibility of interchange the time harmonic electric current densities and the resulting electromagnetic fields in Maxwell's equations for time-invariant linear media\footnote{The inverse, reciprocity, implies that the dependence between an oscillating current and the generated electric field is unchanged if one interchanges the points where the current is placed and where the field is measured. This is why antennas radiation and receiving patterns are similar, and work equally well as transmitters or receivers.}. Systems with nonsymmetric dielectric tensor $\boldsymbol{\epsilon}$ are nonreciprocal systems \cite{potton2004reciprocity}. In this thesis, part of the work is devoted to the study of magnetoplasmonic systems, where reciprocity does not hold.\\
The beginning of magnetoplasmonics dates back to the 70's, when fundamental studies of the effect of an external magnetic field on the dielectric function of the electron plasmas in semiconductors \cite{chiu1972magnetoplasma, brion1972theory, chiu1972magnetoplasma_2, palik1973surface, hartstein1974observation, hartstein1975investigation} where made. In the following two decades the subject remained active due to the possible application in magnetic recording \cite{betzig1992near}. In the last decade, with the possibility of design nanoscale systems almost at will, the interest in magnetoplasmonic systems was renewed, not only by the possible applications in sensing \cite{sepulveda2006highly} or telecommunications \cite{temnov2010active}, but also to obtain fundamental information about the interaction between magnetic field, magneto-optic properties and surface plasmons. Also we should not ignore that this interest is directly connected to the strong development of the plasmonics in recent years, a branch of nanophotonics that aims the understand and control of light using the surface plasmon resonances sustained in a metal nanostructure \cite{maier2007plasmonics, barnes2003surface}.\\
Surface plasmons are electromagnetic waves coupled to plasmons, that are oscillations of the conduction electrons with respect to the fixed ions in a metal, and mathematically appear as strongly localised solutions of Maxwell's equations. They occur in interfaces that separate two media with different signs in the dielectric function, like in a dielectric, $\Re\left({\epsilon_{\textnormal{diel}}}\right)>0$, and a metal, $\Re\left({\epsilon_{\textnormal{met}}}\right)<0$. Surface plasmons exhibit intense resonances when losses are small, $\Im\left({\epsilon_{\textnormal{met}}}\right)<\Re\left({\epsilon_{\textnormal{met}}}\right)$, and are present in a huge variety of metals such as gold or silver, and structures like thin films or nanoparticles \cite{stockman2011nanoplasmonics}. Its high sensitivity to the optical properties of the surrounding dielectric media make them ideal for sensing applications \cite{anker2008biosensing, homola2008surface,cole2006optimized}.\\
 Surface plasmons can be categorised in two types: surface plasmon polaritons (SPP) and localised surface plasmons (LSP). SPP, also called propagating surface plasmons \cite{zayats2003near}, are those sustained by planar interfaces, like in continuous films, and are able to confine the electromagnetic field is nanoscopic volumes, below the diffraction limit. SPP have a propagating character in the parallel direction to the interface, decay exponentially in the direction perpendicular to the surface and their propagation is restricted to a finite distance due to ohmic losses.\\
LSP or particle plasmons are plasmons that are sustained by elements with volume of the order or smaller than the exciting radiation wavelength \cite{pelton2008metal}. Due to their isolated character, these plasmons are not propagating. The LSP excitations are identified by a well defined peak in the optical absorption spectra. The position of these peaks are very sensitive to the optical surrounding environment, as well as shape, size and composition of the nanostructure \cite{link1999spectral,kelly2003optical,aizpurua2003optical,lee2006gold,rodriguez2012gold}.\\
Plasmonic systems have been employed in a vast range of fields, that goes from condensed matter to biology \cite{Stockman:2011ja}. We must highlight the promising applications is cancer treatment \cite{lal2008nanoshell, boisselier2009gold, huang2009gold}, sensing \cite{becker2010optimal,cobley2009shape,stewart2008nanostructured,homola1999surface,teo2010gold} and solar energy conversion \cite{atwater2010plasmonics, schuller2010plasmonics}.\\
An important step forward in plasmonic systems is the control of the plasmon properties using external parameters, called active plasmonics. Some developments have been made in this direction, where electric field \cite{dicken2008electrooptic, dionne2009plasmostor, agrawal2011integrated}, temperature \cite{nikolajsen2004surface, gosciniak2010thermo} or electromagnetic waves \cite{pacifici2007all, pala2008nonvolatile, macdonald2008ultrafast, temnov2009femtosecond} have been used as external agents. An interesting alternative is the use of an external magnetic field \cite{temnov2010active, belotelov2011enhanced, banthi2012high, armelles2013magnetoplasmonics, chin2013nonreciprocal}.\\
The modification of the optical properties by the presence of magnetic fields where initially observed by Faraday \cite{faraday1846experimental} and Kerr \cite{Jkerr1877, Jkerr1878} in the XIX century. Faraday, in transport experiments in glasses, detected a rotation of the polarisation plane that was linearly proportional to the applied magnetic field in the propagation direction.  Kerr, in reflection experiments, observed a rotation of the polarisation plane in a magnetised material. In recent years, magneto-optics have received attention recently due to the possible applications, such as nonreciprocal optical isolators \cite{van2006transverse,dotsch2005applications,espinola2004magneto}, sensors \cite{sepulveda2006highly}, plasmonic interferometry \cite{martin2010enhancement, martin2013plasmonic}, random-lasers \cite{pinheiro2008statistics}, etc. With the goal of increase the magneto-optical response and reduce the size of the devices, magnetic have been combined with plasmonic materials to take advantage of the electromagnetic field enhancement. Also new structures, materials and geometries have been explored \cite{armelles2013magnetoplasmonics}.\\
In this thesis we focus our attention in the modification of the optical properties of complex systems formed by interacting nanoparticles that sustain LSP, when in the presence of a magnetic field.

\section{Thesis structure}

%This thesis consists in two parts, preceded by a general and a technical introduction.\\
%In Part (\ref{part1}) we study light scattering phenomena in systems with few number of scattering elements, but with rich geometry and presenting magneto-optical activity. We study hybrid magnetoplasmonic systems, where we study the importance of interactions between the forming elements. Also we investigate the properties of a emitter-antenna system when the antenna is magneto-optically active.\\
%In Part (\ref{part2}) we study the scattering of light by a large number of elements that are randomly correlated. We investigate the relation between structural correlations with the statistical properties of emission dynamics.\\
%In the end we present the most relevant conclusion gathered from this thesis.

The structure of this thesis is as follows:\\
%In chapter (\ref{chap1}) we make a short overview of light scattering with matter, in particular we focus our attention in complex systems as well as magneto-optical ones.\\
In chapter (\ref{chap:2}) we present the fundamental tools used in the thesis, such as the Green function concept and the fields generated by a point source, the Discrete Dipole Approximation, the polarisability of a small anisotropic particles or the dielectric tensor of a scattering element in the presence of a magnetic field, among others.\\

After the introductory chapters, we present the main research results in two parts, each one with three chapters.\\
In Part (\ref{part1}) we study light scattering phenomena in systems with a few number of scattering elements, but with a rich geometry and presenting magneto-optical activity. \\
In chapter (\ref{chap3}) and chapter (\ref{chap4}) we use a simple, yet comprehensive, model based on the discrete dipole approximation, to study hybrid magnetoplasmonic systems. 
%  where is analysed the influence of a magneto-optical activity over a plasmonic system. 
The considered systems are formed by two nanodisks. In chapter (\ref{chap3}) one of the nanodisks is plasmonic and the other magnetoplasmonic, and we analyse the importance of the interactions between elements of the system and their contribution to the global magneto-optical response. %Due to the coupling and the existence of a magnetoplasmonic element, we study the induced magneto-optical activity in the plasmonic nanodisk by the magnetoplasmonic element.
In chapter (\ref{chap4}) we present a detailed study of the interaction effects when plasmonic and magnetoplasmonic nature of the nanodisk components is changed. \\
In chapter (\ref{chap5}) we study the properties of an emitter-antenna system when the antenna is magneto-optically active. We study two different configurations, where we explore the modifications of the decay rate and the radiated pattern fields with the application of the external magnetic field.\\

In Part (\ref{part2}) we study the statistical properties of light emission in strongly scattering systems with random but correlated structures. \\
In chapter (\ref{chap6}) we study the statistical decay rate dependence in a bi-dimensional strongly scattering system as a function of the particles correlations. We relate the presence of long-range correlations with large fluctuations of the decay rates. We show that not only the parameters governing the statistical distributions of lifetimes are non-universal, but also that the nature of the distributions themselves depend on the details of the system.\\
% as well as short-range correlations with non-Gaussian long-tailed behaviour.
%Here we try to prove that non-universality is richer than expected, because not only the parameters governing the statistical distributions of lifetimes are non-universal, but also the nature of the distributions themselves depend on the details of the system.
In chapter (\ref{chap7}) and chapter (\ref{chap8}) we study the emission decay rate statistics in three-dimensional systems with structural correlations. 
In chapter (\ref{chap7}) we present and analyse the structural model. We consider a cluster of strongly confined particles interacting through a Lennard-Jones potential. We perform a thermodynamical study of the system and, in particular, we focus our attention in the solid-liquid phase transition. In chapter (\ref{chap8}) we consider the same structural system, where we perform light emission calculations at the phase transition. The goal in this chapter is to link the statistical properties of emission dynamics with the thermodynamics of the system, and show that there are fundamental differences between lifetime statistics and transport properties.\\
In chapter (\ref{chap9}) we present the most relevant conclusions of this thesis.

Part of the work developed in this thesis consisted in the implementation of the Discrete Dipole Approximation numerical method for a $\mathcal{N}$ dipole system. The software was designed to make transport and emission calculations, with particles characterised by their corresponding dielectric tensor, shape and volume. Also we implemented the structural program used to generate and analyse clusters.\\
The experimental work, presented in chapter (\ref{chap3}) and chapter (\ref{chap4}), was made in collaboration with \textit{Magnetic Nanostructures and Magneto-Plasmonics} group from Instituto de Microelectr�nica de Madrid.\\

\begin{comment}

\section{Historical overview}

 Although modern studies about interaction of light with matter exhibit a high degree of complexity, the initial questions about light where about its deepest nature.\\
The understanding of the nature of light is intrinsically connected to the human quest for God. From Ancient Egypt to middle ages, religion and belief was behind the efforts for the light comprehension. Greek scholars, arab philosophers and catholic scholars, during 18 centuries, slowly initiates a separation between religion, divinity and God from the scientific interpretation of light's nature. From XIII century, the scientific method was the used technique for the light comprehension. During this period a number of individuals have distinguished themselves by their outstanding works, redefining and changing the knowledge barriers.\\
Empedocles (490 - 430 BC), author of the first comprehensive theory of light and vision, proposed that visual perception where accomplished by rays of light emitted by the eyes, called emission theory. Euclid (mid-4th century BC - mid-3rd century BC) based on Empedocles' work, realises that light travels in straight line and connects mathematics, in particular geometry, with the light theory. Al Hazen (965 - 1040), an arab polymath, refutes the emission theory and proposes the intromission theory, which states that visual perception comes from something representative of the object (later established to be rays of light reflected from it) entering the eyes. Also was the author of several fundamental texts about light and vision, studies about laws of refraction and various physical phenomena such as the rainbow, shadows, eclipses. Al Hazen represents a cornerstone between classical and modern optics, where light leaves the abstract and get into the real world, being converted in laws and mathematical rules.
To the occidental world, the scientific revolution arrives with Roger Bacon (1214 - 1294), an english catholic scholar. Inspired by the legacy of Al Hazen, Bacon studies the physiology of eyesight, the anatomy of the eye and the brain. Describes with scientific rigour the operation of many optical tools, like mirrors and lenses and create the basis for the invention of instruments such as the telescope and the microscope. Bacon, an advocator of the modern scientific method and considered the first occidental scientist, is the biggest responsible for the separation of light from religion.\\
In the XVII century the nature of light was a subject of speculation and research by nearly all great scientists. Snell's law, Newton's ring, Huygens' principle and Fermat's principle date from that time. At same time two theories ruled the scientific stage about the nature of light. The first, supported by Isaac Newton (1642 - 1727), argued that light was composed of particles or corpuscles, which were refracted by accelerating into a denser medium. On the other side was Huygens (1629 - 1695), defending the idea that light was something in the ether as sound is in air. This wave theory of light had difficulties in being accepted principally due to the inexplicable polarisation phenomenon.\\
Only in the XIX century decisive steps where taken. Thomas Young (1773 - 1829) studied diffraction phenomena and justify the obtained pattern of maxima and minima in the shadow space behind a hair was caused by interference of waves, coming from both sides of  hair. Young also suggested that polarisation phenomena was obtained by transverse vibrations of the ether. Augustin-Jean Fresnel (1788 - 1827), using Young's idea of transverse vibration was able to derive these intensity rules theoretically.\\
The climax was reached in the XIX century with James C. Maxwell (1831 - 1879). Maxwell describes, from the experimental laws of Coulomb, Amp�re and Faraday, how electric and magnetic fields are generated and altered by each other and by charges and currents. Also unifies the electric and optical phenomena into the electromagnetic theory. The experimental evidence came 15 years later by the hand of Heinrich Hertz ( 1857 - 1894), where he detected the signal of oscillating charges through an antenna. This crucial moment was accompanied by important developments in mathematics perpetrated by Poisson, Cauchy, Green, Kirchhoff, among others. \\
From this moment, scattering problems start to be faced with more powerful tools. John Strutt (1842 - 1919), also known as Lord Rayleigh, studied the elastic scattering of light by particles much smaller than the wavelength of the light. Of special importance because describes systems whose size is similar to the wavelength, Gustav Mie (1869 - 1957) formulated and presented the full solution of the scattering of light by an homogeneous sphere.\\
This period came to an end with the beginning of the XX with the rise of quantum mechanics. Albert Einstein (1879 - 1955) dismantle the ether hypothesis and shows that light don't need a material medium to propagate. Plank's quantised energy, Einstein explanation of the photoelectric effect, Lewis introduction of the "photon" term, slowly gives form to the new way of understand light. In the end of 1920's Paul Dirac (1902 - 1984), Wolfgang Pauli (1900 - 1958), Victor Weisskopf (1908 - 2002) and P. Jordan (1902 - 1980) applied the fresh quantum mechanics to fields instead of single particles, resulting in the quantum field theories. This area of research culminated with quantum electrodynamics formulation by Richard Feynman (1918 - 1988), Freeman Dyson (1923 -), Julian Schwinger (1918 - 1994) and Shin'ichir\={o} Tomonaga (1906 - 1979).\\
In the last four decades there have been a comeback to the XIX's century classical problems, in particular scattering by small particles and wave propagation in complex media. Developments in astrophysics, chemistry, biology or general condensed matter physics are some of the sources from this new interest. The arise of fabrication technics at nanoscale made nanotechnology a reality. The increase of computing power turned unsolvable into solvable problems.

This brief and oversimplified history of the subject it aims to frame historically the thesis. Interesting approaches to history of optics and science can be found in textbooks such as those by Luis Bernardo \cite{bernardo2005historias, bernardo2007historias, bernardo2010historias} or Geroge Sarton \cite{sarton1968introduction}.


\end{comment}







